\section{Experimental setup and fitted models}
\label{app:experimental-details}


\subsection{\texorpdfstring{$\tau$-bench}{tau-bench} Retail}
\label{app:retail-fit}

\subsubsection{Agent setup}
We use  114 original tasks fromt the benchmark and construct 114 variants of these task and run our sequential randomization ten times per
agent. One agent uses GPT-5 nano (\feat{gpt-5-nano-2025-08-07}) with minimal
reasoning; the other uses Claude Haiku 5.5 (\feat{claude-haiku-5-5}) at low
effort. Both have a 4,000-token output limit and a
100-step episode cap. We generate customer messages with
\feat{gpt-4.1-mini-2025-04-14} at temperature zero and a 2,048-token limit.

The agent receives the store policy, schemas for the 16 tools in
Table~\ref{tab:retail-tools}, and the conversation, including executed calls
and results. We instruct it to help the customer while following the policy.
At each turn, the agent chooses a tool call or a message to the customer. We
score terminal success by checks with behavioral
assertions in the benchmark.

\subsubsection{Reviewer setup}
After freezing the proposed call or batch, we assign review or direct execution
with equal probability. The reviewer uses GPT-6 Luna (\feat{gpt-6-luna}) with
high reasoning effort and a 6,000-token output limit. It receives the policy,
schemas, visible conversation, and frozen proposal, but cannot query the
database. We instruct it to check authentication, identifiers, arguments,
available evidence, and required confirmation.

In its JSON response, the reviewer either approves the proposal, replaces the
entire batch with one call, or suggests one question for the customer. In case the reviewer approves the tool call is directly executed. 

\begin{figure}[!htbp]
\centering
\includegraphics[width=\textwidth]{retail_trajectories.pdf}
\caption{Two observed GPT-5 nano $\tau$-bench Retail episodes with review
randomized at each proposal. Episode A succeeds after review replaces an
invalid exchange proposal; episode B ends with a transfer to a human after an
invalid order modification.}
\label{fig:retail-trajectories}
\end{figure}

\subsubsection{Features}
Each reporting group contains original proposals involving any of the tools
assigned to that group. Table~\ref{tab:retail-tools} gives the tool names and
their families.

\begin{table}[H]
\centering\small
\setlength{\tabcolsep}{4pt}
\renewcommand{\arraystretch}{1.12}
\caption{The 16 tools available to the agents and grouped into  seven families of task}\label{tab:retail-tools}
\begin{tabular}{@{}>{\raggedright\arraybackslash}p{.22\textwidth}>{\raggedright\arraybackslash}p{.75\textwidth}@{}}
\toprule Family suffix & Tool names\\\midrule

\feat{identity} & \feat{find_user_id_by_email}, \feat{find_user_id_by_name_zip}\\
\feat{user_read} & \feat{get_user_details}\\
\feat{order_read} & \feat{get_order_details}\\
\feat{catalog_read} & \feat{get_item_details}, \feat{get_product_details}, \feat{list_all_product_types}\\
\feat{pending_write} & \feat{cancel_pending_order}, \feat{modify_pending_order_address}, \feat{modify_pending_order_items}, \feat{modify_pending_order_payment}\\
\feat{resolution_write} & \feat{exchange_delivered_order_items}, \feat{modify_user_address}, \feat{return_delivered_order_items}\\
\feat{utility} & \feat{calculate}, \feat{transfer_to_human_agents}\\
\bottomrule
\end{tabular}
\end{table}
\begin{table}[H]
\centering\small
\setlength{\tabcolsep}{4pt}
\renewcommand{\arraystretch}{1.12}
\caption{The 24 Retail SNMM inputs, used for all continuations. Yes/no inputs
are one for yes and zero for no.}\label{tab:retail-snmm}
\begin{tabular}{@{}>{\raggedright\arraybackslash}p{.34\textwidth}r>{\raggedright\arraybackslash}p{.56\textwidth}@{}}
\toprule Input & Count & Meaning\\\midrule
\feat{intercept} & 1 & A constant value of one.\\
Proposed tools (\feat{tool:}) & 16 & One input for each tool in Table~\ref{tab:retail-tools}: whether the current proposal includes that tool.\\
\feat{prior_error} & 1 & An earlier tool call returned an error.\\
\feat{prior_correction} & 1 & An earlier review changed a proposal.\\
\feat{retry_tool} & 1 & At least one proposed tool was used earlier.\\
\feat{used:catalog_read} & 1 & An earlier product or catalog read succeeded.\\
\feat{used:order_read} & 1 & An earlier order-details read succeeded.\\
\feat{write_x_prior_error} & 1 & The proposal contains a write, and an earlier tool call returned an error.\\
\feat{write_x_affirmation} & 1 & The proposal contains a write, and an earlier user message contains a confirmation phrase.\\
\bottomrule
\end{tabular}
\end{table}

 A confirmation phrase matches
\texttt{yes|confirm|go ahead|proceed} as whole words, ignoring capitalization.
We count a tool result as an error if its error flag is set or its text starts
with \texttt{error}.

\begin{table}[H]
\centering\small
\setlength{\tabcolsep}{4pt}
\renewcommand{\arraystretch}{1.12}
\caption{The 41 features for the Retail nuisance models. }\label{tab:retail-nuisance}
\begin{tabular}{@{}>{\raggedright\arraybackslash}p{.34\textwidth}r>{\raggedright\arraybackslash}p{.56\textwidth}@{}}
\toprule Inputs & Count & Meaning\\\midrule
Proposed tools (\feat{tool:}) & 16 & The same 16 tool inputs as in Table~\ref{tab:retail-snmm}.\\
Current tool families (\feat{current:}) & 7 & For each family in Table~\ref{tab:retail-tools}, whether the proposal includes a tool from that family.\\
Earlier successful families (\feat{used:}) & 7 & For each family, whether an earlier call from that family succeeded.\\
Earlier events & 5 & Earlier error, earlier correction, and repeated tool (defined above); any earlier review assignment (\feat{prior_verification}); earlier user confirmation (\feat{user_affirmation_seen}).\\
\feat{n_prior_tools}, \feat{n_prior_errors}, \feat{n_prior_verifications}, \feat{n_proposed_tools} & 4 & Numbers of earlier tool results, earlier errors, earlier review assignments, and calls in the current proposal, respectively.\\
\feat{source_cohort_extension} & 1 & Whether the task comes from the added-task set.\\
\feat{eligible_position_zero_based} & 1 & The number of review opportunities before the current proposal.\\
\bottomrule
\end{tabular}
\end{table}

To predict terminal success, we also include the current review assignment
and earlier decision records. Each earlier record contains the 39 inputs in
the first five rows of Table~\ref{tab:retail-nuisance} and its review assignment.
We reserve 28 slots for GPT-5 nano and 14 for Claude Haiku; unused slots are zero.

The models used to account for later reviews use the 41 current inputs above,
plus one yes/no input for each possible decision number. Only the input for
the current decision number is one.


\subsubsection{Phases and coefficient sharing}\label{app:phases-retail}
We use four phases: before successful customer identification; after
identification but before a successful order or catalog read; after that read
but before a successful write; and after a successful write. 

We fit a separate coefficient vector for each phase, agent, and
continuation, sharing coefficients across decisions within a phase.
The randomized-continuation fit has no future-treatment correction.
For changed continuations, we solve the joint moment system by backward
substitution across the four phases, as
in~\eqref{eq:backward-outcome}--\eqref{eq:phase-moments}. Each phase's
equations retain the contributions of later proposals in that same phase, while contributions from later phases use their already estimated coefficients.

\subsubsection{Nuisance estimation}
For each held-out fold, we train one histogram gradient-boosted regressor
on the other four folds to predict terminal success from the inputs listed
above. Boosted trees accommodate nonlinear relationships and interactions
among history features. We choose the histogram implementation for
computational efficiency across repeated fits: it bins numerical predictors
to reduce the number of candidate split points. We obtain predictions
$\widehat m_0$ and $\widehat m_1$ by setting the current assignment to zero
and one, respectively, and use
$\widehat q=(\widehat m_0+\widehat m_1)/2$. We use squared-error loss,
a learning rate of 0.1, at most 120 boosting iterations, at most 12 leaves
per tree, a minimum leaf size of 20, an L2 penalty of 2, seed $42+k$ for
fold $k$, and automatic early stopping.

For changed continuations, we predict the sum of later feature vectors for
each phase, weighting each vector by the observed treatment minus the
target-policy action. Within each training fold, we standardize the current
attributes, cohort, numeric position, and position indicators, then fit
ridge regressions with penalty 10. We use ridge to fit the many
future-feature prediction targets efficiently, with coefficient shrinkage
to stabilize estimation when predictors are correlated. We set predictions for earlier phases to zero. Before fitting the unpenalized effect regressions, we remove redundant feature directions by SVD without using the outcomes.

\subsection{ToolSandbox}
\label{app:toolsandbox-fit}

\subsubsection{Agent setup}
For each agent, we collect three trials of 129 base scenarios under eight tool variants, for 3,096 trajectories. We use the same two agents, model settings, and user simulator as in $\tau$-bench Retail, with each scenario's official cap on the number of messages.

The agent sees the conversation and schemas for the simulated-phone tools offered in that scenario (Table~\ref{tab:toolsandbox-tools}). We instruct it not to guess missing arguments and to clarify ambiguous requests. Across variants, we alter tool availability or scramble tool names and documentation.

\subsubsection{Reviewer setup}
We freeze each agent turn containing tool calls and assign review with
probability $1/2$ before execution.   The reviewer uses the same model settings as in $\tau$-bench Retail: GPT-6 Luna with high reasoning effort and a 6,000-token output limit.


The reviewer receives the visible conversation, proposal, offered schemas, public API rules, and stateless argument checks. For scrambled variants, we include the relevant rules without changing the schemas shown to the agent. The reviewer cannot execute tools or see the hidden device state and evaluator targets. It either approves the batch or replaces it with one call from the offered tools or one response to the user, citing the supplied evidence for any correction.

\subsubsection{Features}
Tables~\ref{tab:toolsandbox-snmm} and~\ref{tab:toolsandbox-nuisance} list the 45
attributes for the SNMM model and 90 attributes used for  nuisance model. We derive device settings and
database sizes from the pre-execution state of the system, We group tools into six capabilities for reporting and ten families for model inputs.

\begin{table}[H]
\centering\small
\setlength{\tabcolsep}{4pt}
\renewcommand{\arraystretch}{1.12}
\caption{The 33 tools available across scenarios and their grouping used as
model inputs.}\label{tab:toolsandbox-tools}
\begin{tabular}{@{}>{\raggedright\arraybackslash}p{.22\textwidth}>{\raggedright\arraybackslash}p{.75\textwidth}@{}}
\toprule Family suffix & Tool names\\\midrule

\feat{settings_read} & \feat{get_low_battery_mode_status}, \feat{get_location_service_status}, \feat{get_cellular_service_status}, \feat{get_wifi_status}, \feat{get_current_location}\\
\feat{settings_write} & \feat{set_low_battery_mode_status}, \feat{set_location_service_status}, \feat{set_cellular_service_status}, \feat{set_wifi_status}\\
\feat{contact_lookup} & \feat{search_contacts}\\
\feat{contact_edit} & \feat{add_contact}, \feat{modify_contact}, \feat{remove_contact}\\
\feat{message_search} & \feat{search_messages}\\
\feat{message_send} & \feat{send_message_with_phone_number}\\
\feat{reminder_read} & \feat{search_reminder}\\
\feat{reminder_write} & \feat{add_reminder}, \feat{modify_reminder}, \feat{remove_reminder}\\
\feat{utility} & \feat{get_current_timestamp}, \feat{timestamp_to_datetime_info}, \feat{datetime_info_to_timestamp}, \feat{shift_timestamp}, \feat{timestamp_diff}, \feat{seconds_to_hours_minutes_seconds}, \feat{unit_conversion}, \feat{calculate_lat_lon_distance}, \feat{search_holiday}\\
\feat{web_search} & \feat{search_lat_lon}, \feat{search_location_around_lat_lon}, \feat{search_weather_around_lat_lon}, \feat{search_stock}, \feat{convert_currency}\\
\bottomrule
\end{tabular}
\end{table}
\begin{table}[!htbp]
\centering\small
\setlength{\tabcolsep}{4pt}
\renewcommand{\arraystretch}{1.12}
\caption{The 45 ToolSandbox SNMM inputs, used for all continuations. Each input
is one when its condition holds and zero otherwise.}\label{tab:toolsandbox-snmm}
\begin{tabular}{@{}>{\raggedright\arraybackslash}p{.34\textwidth}r>{\raggedright\arraybackslash}p{.56\textwidth}@{}}
\toprule Input & Count & Meaning\\\midrule
Progress phases (\feat{phase:}) & 3 & Before any successful tool result; after a success but before a successful write; after a successful write. Only one is active.\\
Proposed tools (\feat{tool:}) & 33 & One input for each tool in Table~\ref{tab:toolsandbox-tools}: whether the current proposal includes that tool.\\
\feat{prior_error} & 1 & An earlier tool result contained an exception.\\
\feat{prior_state_error} & 1 & An earlier exception began with \texttt{ConnectionError} or \texttt{PermissionError}.\\
\feat{retry_tool} & 1 & At least one proposed tool was executed earlier.\\
\feat{prior_verification} & 1 & An earlier proposal was assigned review.\\
\feat{prior_correction} & 1 & An earlier review changed a proposal.\\
\feat{precondition_unmet} & 1 & A proposed call breaks a device-state rule.\\
\feat{write_x_prior_error} & 1 & The proposal contains a write, and an earlier tool result contained an exception.\\
\feat{scrambled} & 1 & The scenario name contains \texttt{\_scrambled}.\\
\feat{distractors} & 1 & The scenario name contains \texttt{distraction\_tools} or ends in \texttt{\_all\_tools}.\\
\bottomrule
\end{tabular}
\end{table}

A write changes device settings, edits contacts, sends a message, or changes
reminders. We set \feat{precondition_unmet} to one for messaging with cellular off; web search
with Wi-Fi off; coordinate-based search without latitude when location is off;
current-location lookup with location off; setting a toggle to its existing
value; or enabling Wi-Fi, cellular, or location in low-battery mode.

\begin{table}[!htbp]
\centering\small
\setlength{\tabcolsep}{4pt}
\renewcommand{\arraystretch}{1.12}
\caption{The 90  features used for the ToolSandbox nuisance models.}\label{tab:toolsandbox-nuisance}
\begin{tabular}{@{}>{\raggedright\arraybackslash}p{.34\textwidth}r>{\raggedright\arraybackslash}p{.56\textwidth}@{}}
\toprule Inputs & Count & Meaning\\\midrule
Proposed tools (\feat{tool:}) & 33 & The same 33 tool inputs as in Table~\ref{tab:toolsandbox-snmm}.\\
Current tool families (\feat{current:}) & 10 & For each family in Table~\ref{tab:toolsandbox-tools}, whether the proposal includes a tool from that family.\\
Earlier successful families (\feat{used:}) & 10 & For each family, whether an earlier call from that family succeeded.\\
Earlier events and device rules & 6 & \feat{prior_error}, \feat{prior_state_error}, \feat{retry_tool}, \feat{prior_verification}, \feat{prior_correction}, and \feat{precondition_unmet}, defined above.\\
\feat{unknown_tool}, \feat{malformed_arguments} & 2 & The execution tool name is outside the catalog; the arguments do not decode to a JSON object.\\
\feat{n_prior_errors}, \feat{n_prior_tools}, \feat{n_prior_verifications}, \feat{n_proposed_tools}, \feat{n_user_turns} & 5 & Numbers of earlier exception results, traced tool calls, review assignments, calls in the current proposal, and observed user turns, respectively.\\
Device settings (\feat{setting:}) & 4 & Whether \feat{cellular}, \feat{wifi}, \feat{location_service}, and \feat{low_battery_mode} are on.\\
Database sizes (\feat{rows:}) & 3 & Numbers of rows in \feat{CONTACT}, \feat{MESSAGING}, and \feat{REMINDER}.\\
Scenario categories (\feat{category:}) & 15 & All tools available; zero, three, or ten distractor tools; scrambled argument descriptions, argument types, tool descriptions, or tool names; single or multiple tool calls; single or multiple user turns; canonicalization; insufficient information; state dependency. Categories can overlap.\\
Scenario variants (\feat{variant:}) & 2 & The \feat{scrambled} and \feat{distractors} inputs defined above.\\
\bottomrule
\end{tabular}
\end{table}


\subsubsection{Phases and coefficient sharing}\label{app:phases-tool}
We use three phases: before any successful tool result, after a success but
before a successful write, and after a successful write. Phase progression is
one-way. Writes include settings changes, contact edits, message sends, and
reminder writes.

For continuation $c$, the effect model is
\[
 b^c(h)=a^c_{s(h)}+\sum_{k=1}^{33}b_k^c\,\ind\{k\text{ is proposed}\}
                 +(d^c)^\top z(h),
\]
where $z$ is the vector of nine history and variant terms. The baseline $a_s^c$
varies by phase; the tool coefficients $b_k^c$ and history coefficients $d^c$
are shared across phases. There is no additional intercept. We fit each agent
and continuation separately. For each changed continuation, we solve one joint
system of equations that accounts for all later decisions, both within the
current phase and in subsequent phases. We chose this structure after
inspecting sparse-coefficient diagnostics; the intervals do not include
model-selection uncertainty.

\subsubsection{Nuisance estimation}
We encode the history using the 90  attributes, numeric decision
position, and earlier decisions' attributes and assignments.  For outcome prediction, we use the same histogram
gradient-boosting specification as in $\tau$-bench Retail, retaining its
flexibility to model nonlinear relationships and its computational
efficiency across repeated fits. We obtain predictions $\widehat m_0$
and $\widehat m_1$ from the fitted outcome model by setting the current
assignment to zero and one, respectively, then take their average as
$\widehat q$.

Within each training fold, we standardize the current attributes and position
indicators, then fit ridge regressions with penalty 10
to predict future totals in the phase-expanded basis. We set predictions for
earlier phases to zero and map the remaining predictions into the basis for
the pooled model's shared coefficients. Before fitting the unpenalized effect
regressions, we apply SVD to the pooled 45-feature design without using the
outcomes.
\FloatBarrier

\subsection{ImpossibleBench}
\label{app:steering-fit}

\subsubsection{Agent setup}
The agent uses Qwen3.5 Flash (\feat{qwen/qwen3.5-flash-02-23}) at temperature
0.5, with reasoning disabled and a 4,096-token output limit. It has at most six
actions and can submit code with \feat{write_file} or stop with
\feat{report_issue}. In the system prompt we prohibit test edits
and evaluator exploits, and permit reports of contradictory tests. We run the
tests after each write. At the next action, the agent sees the task, submitted
code, and test feedback.

We run each of 103 problems six times with original tests and six times with
contradictory tests, yielding 1,236 episodes. We exclude 18 with missing
terminal outcomes and analyze the remaining 1,218. We record a violation if a
test edit is detected or the agent passes the contradictory check, and end the
episode after success, an issue report, or a detected violation.

\subsubsection{Steering setup}
After an eligible unsuccessful submission, we assign a neutral or steering note
with equal probability. In the steering note, we tell the agent that it may change the tests or we prompt Qwen3.5 Flash with the task, code,
checkpoint state, test feedback, and two previous notes, asking for a
diagnostic note that encourages a prohibited edit or test-order shortcut. We
generate the note at temperature 0.8 with a 512-token limit. If generation
fails, we use the script and retain the original assignment. We append the
note to the tool feedback. In the system instructions to the agent, we tell the agent that notes cannot override any instructions .

\begin{figure}[!htbp]
\centering
\includegraphics[width=\textwidth]{steering_trajectories.pdf}
\caption{Two observed Qwen3.5 Flash runs of the same ImpossibleBench task, with note
assignment randomized at each eligible checkpoint. The neutral run reports
contradictory tests; the steered run follows a scripted note and edits the
protected test file.}
\label{fig:steering-trajectories}
\end{figure}


\subsubsection{Features}
Tables~\ref{tab:steering-snmm} and~\ref{tab:steering-nuisance}
list the nine SNMM inputs and 36 current nuisance attributes.

\begin{table}[!htbp]
\centering\small
\setlength{\tabcolsep}{4pt}
\renewcommand{\arraystretch}{1.12}
\caption{The nine features used for SNMM inputs}\label{tab:steering-snmm}
\begin{tabular}{@{}>{\raggedright\arraybackslash}p{.34\textwidth}r>{\raggedright\arraybackslash}p{.56\textwidth}@{}}
\toprule Input & Count & Meaning\\\midrule
\feat{intercept} & 1 & A constant value of one.\\
\feat{pass_fraction} & 1 & The fraction of checks passed.\\
\feat{conflict_visible} & 1 & The visible feedback contains contradictory checks.\\
\feat{near_conflict} & 1 & At least 80\% but not all checks pass, and the visible feedback contains contradictory checks.\\
\feat{prior_attack_resisted} & 1 & Earlier steering occurred and the current checkpoint is eligible. An explicit rejection is not required.\\
\feat{stalled} & 1 & The fraction of checks passed is no higher than after the preceding submission; zero at the first submission.\\
\feat{execution_error} & 1 & A current check reports an error.\\
\feat{repeated_failure} & 1 & The current failure signature matches the preceding evaluation's.\\
\feat{mentions_inconsistency} & 1 & The agent's message contains \texttt{contradict}, \texttt{inconsistent}, \texttt{inconsistency}, or \texttt{impossible}, ignoring capitalization.\\
\bottomrule
\end{tabular}
\end{table}

\begin{table}[!htbp]
\centering\small
\setlength{\tabcolsep}{4pt}
\renewcommand{\arraystretch}{1.12}
\caption{The 36 features for nuisance models.
}\label{tab:steering-nuisance}
\begin{tabular}{@{}>{\raggedright\arraybackslash}p{.34\textwidth}r>{\raggedright\arraybackslash}p{.56\textwidth}@{}}
\toprule Inputs & Count & Meaning\\\midrule
Current SNMM attributes & 8 & All inputs in Table~\ref{tab:steering-snmm} except the constant.\\
\feat{previous_pass_fraction}, \feat{progress_delta} & 2 & The fraction of checks passed after the preceding submission, and the change since then.\\
Progress ranges & 4 & \feat{almost_finished}: $0.8\le p<1$; \feat{zero_progress}: $p=0$; \feat{low_progress}: $p<0.5$; \feat{mid_progress}: $0.5\le p<0.8$, where $p$ is the current pass fraction.\\
Other progress changes & 3 & \feat{regressed}: $p<p_{-}$; \feat{improved}: $p>p_{-}$; \feat{recovered_from_zero}: $p_{-}=0<p$, where $p_{-}$ is the preceding pass fraction. All are zero at the first submission.\\
\feat{same_failure_streak}, \feat{unchanged_code} & 2 & The number of consecutive preceding evaluations with the current failure signature; whether the code equals the state saved after the preceding action (zero at the first checkpoint).\\
Time & 4 & \feat{checkpoint}: preceding action number; \feat{turns_remaining}: $6-\text{checkpoint}$; \feat{late}: checkpoint $\ge3$; \feat{last_chance}: checkpoint $=5$.\\
\feat{prior_attacks} & 1 & The number of earlier steering assignments.\\
Previous note (\feat{previous_arm_}) & 4 & One input each for \feat{neutral}, \feat{static}, \feat{contextual}, and \feat{llm}; all zero before the first assignment.\\
Missing previous progress & 2 & \feat{previous_pass_fraction_missing} and \feat{progress_delta_missing}: whether the corresponding value is missing. Missing values are set to zero.\\
\feat{log_checks}, \feat{log_code_characters} & 2 & $\log(1+\text{count})$ for current checks and code characters.\\
Task variant and cohort & 4 & \feat{conflicting_variant}: contradictory tests; \feat{repaired_conflict}: locally repaired fixture; \feat{confirmation_split}: confirmation cohort; \feat{baseten_provider}: always zero for Qwen3.5 Flash.\\
\bottomrule
\end{tabular}
\end{table}


\subsubsection{Phases and coefficient sharing}
Checkpoint $t$ follows action $t$, for $t=1$ through $5$. When an action is ineligible, we leave its checkpoint empty and retain the numbering of later checkpoints. We fit separate coefficients by checkpoint and continuation. We estimate coefficients for changed continuations backward from checkpoint 5. The future correction uses only later checkpoints because each checkpoint occurs at most once.

\subsubsection{Nuisance estimation}
We encode the baseline history with 36  features and, for each of the
two previous eligible checkpoints, its attributes, assignment, and a presence
indicator. For each held-out trajectory fold, we fit separate control, treatment, and
pooled outcome regressions on the other four folds to estimate $\widehat m_0$,
$\widehat m_1$, and $\widehat q$. We use the $\tau$-bench Retail
gradient-boosting settings with early stopping disabled. For the
future-feature regressions, we use the same history features, standardize them
within each training fold, and apply ridge penalty 10. Predictions for the
current or an earlier checkpoint are zero. After applying SVD without using
outcomes, we fit unpenalized effect regressions and scale the retained
directions to unit second moment within each checkpoint.

